\subsection{Minimality of a smooth Hopf normalization}

\begin{statement}{P4.24}{Proved}
In the setting of Part~4.2, let $S$ be an irreducible
singular fibre support with smooth normalization $T$.
If the minimal model of $T$ is a Hopf surface, then
$T$ is already minimal. The reduced conductor on $T$
and its reduced image in $S$ are disjoint unions of
smooth elliptic curves. Each map between their
components is a finite unramified covering, and
$e(S)=0$.
\end{statement}
\begin{proof}
Write $\nu:T\to S\hookrightarrow Y$ for the finite
holomorphic normalization map followed by the closed
embedding. We use the effective integral conductor
divisor $D$ with $K_T+D$ torsion. Let $T\to T_0$
be the sequence of $k$ point blowups giving the
minimal Hopf model. At each step the conductor
coefficient on the exceptional curve is $\mu-1$,
where $\mu$ is the multiplicity of the conductor
at the centre. Effectiveness therefore puts every
centre on the current conductor.

There is a finite unramified primary Hopf cover
$T'_0\to T_0$. Pulling back every blowup gives a
finite unramified cover $T'\to T$ and a blowup
presentation of $T'$ over $T'_0$. A centre has
exactly the degree of the covering many lifts,
so the new number of blowups is that degree times
$k$. The composite $T'\to Y$ is finite and is
immersive away from the pulled-back conductor.
The rational Chern classes of its pulled-back
ambient tangent bundle vanish. Indeed $b_2(Y)=0$
gives $c_1(T_Y)=c_2(T_Y)=0$ in rational cohomology,
using Poincar\'e duality in degree four. The
argument below applies to this finite map, without
requiring it to be a normalization. Proving
minimality of $T'$ proves $k=0$. We therefore use
$T,T_0,\nu,k$ for the covering objects during the
Chern calculation.

We first construct on $T_0$ a nonsingular
holomorphic foliation tangent to every compact
curve. Khaled, \emph{Real structures on primary
Hopf surfaces}, Theorem~2.1(1) and formula~(1),
gives
\[
 T_0=(\mathbb C^2\setminus\{0\})/\langle F\rangle,
 \qquad F(z,w)=(\alpha z+\lambda w^n,\beta w),
\]
where $0<|\alpha|\le|\beta|<1$ and
$\lambda(\alpha-\beta^n)=0$. We verify the needed
curve statement from these coordinates. Every
compact irreducible curve is smooth elliptic:
adjunction and $b_2(T_0)=0$ give arithmetic genus
one, and a rational normalization would lift to
$\mathbb C^2\setminus\{0\}$ and have constant
coordinate functions.

The map from the fundamental group of such an
elliptic curve to $\pi_1(T_0)=\mathbb Z$ has
nonzero image, since zero image would give the
same impossible compact lift. Its full inverse
image has finitely many connected components.
The Remmert--Stein extension across the origin
is an analytic curve germ; a power $F^h$ fixes
one of its branches. On the normalization of
that branch the induced contraction can be
written $t\mapsto\eta t$, $0<|\eta|<1$.
For completeness, if the initial coordinate
has return map $g(t)=\eta t+O(t^2)$, choose
$|\eta|<\rho<\sqrt{|\eta|}$. On a sufficiently
small disc, $|g(t)|\le\rho|t|$, and the
successive differences in $g^j(t)/\eta^j$
are bounded by a constant times
$(\rho^2/|\eta|)^j|t|^2$. Their convergent sum
gives the asserted linearizing coordinate.

When $\lambda=0$, comparison of coefficients in
$z(\eta t)=\alpha^h z(t)$ and
$w(\eta t)=\beta^h w(t)$ shows that each
nonzero coordinate is a single monomial.
Thus every branch away from the coordinate axes
has parametrization $(a t^p,b t^q)$. Its ratio
$p/q=\log|\alpha|/\log|\beta|$ is independent
of the branch. If such a branch exists, take
coprime positive integers $p,q$ in this ratio
and the vector field
$p z\partial_z+q w\partial_w$.
If no such branch exists, use
$z\partial_z+w\partial_w$.
Both are invariant under $F$, have no zero on
$\mathbb C^2\setminus\{0\}$, and are tangent
to every compact curve.

When $\lambda\ne0$, we have $\alpha=\beta^n$ and
\[
 F^h(z,w)=
 (\alpha^h z+h\lambda\alpha^{h-1}w^n,\beta^h w).
\]
If a normalized invariant branch had
$w(t)=b t^q\ne0$, the coefficient of $t^{nq}$
in the first coordinate identity would give
$0=h\lambda\alpha^{h-1}b^n$, a contradiction.
Thus the only curve is the image of $w=0$.
The vector field $n z\partial_z+w\partial_w$
is invariant, nowhere zero, and tangent to it.
This proves the foliation assertion.

Pull this foliation through the point blowups
and saturate its tangent sheaf. It is a line
bundle $L\subset T_T$. Every strict transform
of a conductor component and every exceptional
curve is invariant. The only singularities of
the foliation are simple saddles. To verify
this assertion and its multiplicities, at a
regular point use coordinates $(x,y)$ with
leaves $y=\text{constant}$. In the chart
$x=u,y=uv$ the saturated generator is
$u\partial_u-v\partial_v$; there is one
saddle, and the other blowup chart is regular
away from it. At an existing saddle use
$b x\partial_x-a y\partial_y$, with positive
integers $a,b$. The two charts have generators
\[
 b u\partial_u-(a+b)v\partial_v,\qquad
 (a+b)u\partial_u-a v\partial_v.
\]
One saddle is replaced by two. Each has
zero ideal $(u,v)$ and length one.
Consequently there are exactly $k$ saddles.

Let $V$ be the divisor measuring the
divisorial zero of the pulled-back transverse
one-form before saturation. Numerically
\[
 c_1(L)=-K_T+V.
\]
The omitted pullback line bundle is topologically
trivial because it comes from $T_0$.
The preceding charts give, at every blowup $\mu$,
\[
 K'=\mu^*K+E,\qquad
 V'=\mu^*V+\varepsilon E,\qquad
 \varepsilon=
 \begin{cases}0&\text{at a regular point},\\
 1&\text{at a saddle}.
 \end{cases}
\]
Starting with $V=0$ and $K_{T_0}^2=0$, we obtain
\[
 K_T^2=-k,\qquad K_TV-V^2=0.
\]
The second identity is preserved because its
change is $\varepsilon^2-\varepsilon=0$.

Put $B=\nu^*T_Y$ and compose the holomorphic
maps $L\to T_T\to B$. This map has no
divisorial zero. Away from the conductor this
follows from immersivity. At the generic point
of a conductor component the foliation is
tangent to that curve, and finiteness of
$\nu$ makes its differential along the curve
nonzero. A nonconstant holomorphic map of
curves has nonzero derivative at the generic
point. Thus the cokernel $Q=B/L$ is a
torsionfree sheaf of rank two. Its double
dual $Q^{**}$ is locally free, and
\[
 z_0=\operatorname{length}(Q^{**}/Q)
\]
is finite. At a point, write the map $L\to B$
as a column $(f_1,f_2,f_3)$ in trivializations.
If this column vanishes, the ideal
$(f_1,f_2,f_3)$ has finite colength. Dualizing
$0\to\mathcal O\to\mathcal O^3\to Q\to0$
gives $\mathcal E xt^1(Q,\mathcal O)
=\mathcal O/(f_1,f_2,f_3)$.
Dualizing $0\to Q\to Q^{**}\to Q^{**}/Q\to0$
and using the two-dimensional regular local
ring identifies this with
$\mathcal E xt^2(Q^{**}/Q,\mathcal O)$,
which has the same length. At each foliation
saddle all three coefficients vanish.
It follows that
\[
 z_0\ge k.
\]

Let $E_B\subset B$ be the saturation of the
generic rank-two image of $d\nu:T_T\to B$.
Its determinant differs from that of $T_T$
by an effective integral divisor $R$, the
divisorial rank-drop divisor:
\[
 c_1(E_B)=-K_T+R.
\]
The support of $R$ is contained in the
conductor. Saturate $E_B/L$ in $Q^{**}$.
The resulting rank-one reflexive subsheaf
is a line bundle $J$, with
$c_1(J)=R-V$. Its torsionfree rank-one
quotient has the form
$\mathcal I_{Z_1}\otimes J_1$ with
$c_1(J_1)=K_T-R$ and
$z_1=\operatorname{length}(Z_1)\ge0$.
All these assertions concern sheaves on the
smooth surface $T$; determinants are computed
off finite sets and therefore determine their
line-bundle classes.

The identities $c_1(B)=c_2(B)=0$ give
\[
 c_2(Q)=(-K_T+V)^2,\qquad
 c_2(Q^{**})=(-K_T+V)^2-z_0.
\]
The exact sequence with $J$ gives another
expression for the same number:
\[
 (-K_T+V)^2-z_0=(R-V)(K_T-R)+z_1.
\]
Using $K_T^2=-k$ and $K_TV-V^2=0$, this becomes
\[
 G(R):=R^2-K_TR-VR=k+z_0+z_1\ge2k.
\]

We compare this expression with the actual
blowup sequence. Push $R$ and $V$ down at
each step; $R$ has integral coefficients.
Writing $R'=\mu^*R+cE$ gives $c\in\mathbb Z$,
without any positivity requirement on $c$.
The intersection calculation is
\[
 G(R')-G(R)=-c^2+(1+\varepsilon)c.
\]
For $\varepsilon=0$ this is at most zero
for every integer $c$; for $\varepsilon=1$
it is $1-(c-1)^2\le1$. On $T_0$ every
divisor is numerically zero, so $G(R_0)=0$.
Consequently
\[
 G(R)\le\#\{\text{blowups at saddles}\}\le k.
\]
The lower bound $G(R)\ge2k$ forces
$k=z_0=z_1=0$. Returning through the finite
unramified cover proves that the original
normalization $T$ is minimal.

It remains to identify the conductor maps,
rather than only their Euler number.
The reduced conductor on a minimal Hopf
surface consists of disjoint smooth elliptic
curves. Since $z_0=0$ on the primary cover,
the differential of $\nu$ along each such
curve is nowhere zero, also before that
cover. Thus every branch of its image is
smooth. At an image point, two different
curve germs would have to come from two
different normalization points. Their
surface branches in the smooth threefold
intersect in a curve germ. This germ is
contained in the conductor image of both
branches; the unique conductor germ at
each normalization point forces the two
image germs to be equal. Therefore the
image conductor is locally a single smooth
curve. Its components are disjoint, and
each conductor map is finite unramified.
Riemann--Hurwitz makes every image component
elliptic. Finally, normalization is an
isomorphism off these curves, so
\[
 e(S)=e(T)-e(D_{\mathrm{red}})
                 +e(\operatorname{Sing}S)=0.
\]
\end{proof}
